% ---------------------------------------------------------------------
% Cover page, matched to the website book-panel card (dark navy panel,
% edition tag, large title with a teal accent word, author line).
% Injected by pandoc as include-before-body.
% ---------------------------------------------------------------------
\begin{titlepage}
\thispagestyle{empty}
\centering
\vspace*{1.2cm}

\begin{tcolorbox}[colback=spaceNavy, colframe=spaceNavy, arc=8pt,
  boxrule=0pt, left=26pt, right=26pt, top=30pt, bottom=30pt,
  enhanced, borderline south={4pt}{0pt}{brandTeal}]
\raggedright
{\sffamily\bfseries\footnotesize\color{white!75!brandTeal}%
 \fcolorbox{white!25!spaceNavy}{white!12!spaceNavy}{\strut~EDITION I\,\textbullet\,ENGINEERING STUDY\,\textbullet\,2026~}}\\[26pt]
{\sffamily\bfseries\fontsize{34}{38}\selectfont\color{white}%
 Space-Based\\ AI Data Centers}\\[10pt]
{\sffamily\itshape\Large\color{brandTeal!75!white}%
 Physics, Mathematics, Simulation, and Economics\\ of an Orbital Compute Constellation}\\[30pt]
{\sffamily\normalsize\color{white!80!spaceNavy}%
 A primary-source engineering study of the orbital data-center concept,\\
 from orbital mechanics and thermal control to debris survivability,\\
 launch economics, and a fully open simulation model.}
\end{tcolorbox}

\vfill

{\sffamily\Large\bfseries\color{brandSlate} Samarjith Biswas, PhD}\\[4pt]
{\sffamily\normalsize\color{brandGray} Research Scientist III}\\[18pt]
{\sffamily\small\color{brandGray}%
 Reference architecture: Google Project Suncatcher (arXiv:2511.19468)\\[2pt]
 Open repository: \texttt{github.com/Samarjithbiswas/space-based-ai-datacenter}\\[2pt]
 Open access \textbullet{} MIT (code) and CC-BY (text)}

\vspace*{1.0cm}
\end{titlepage}
\cleardoublepage
